\paragraph{Human subject annotation protocol.}
Human evaluations (\S\ref{sec:applied}) involved three expert annotators labelling \HumN{} published comparative claims, adjudicated by a senior researcher across \HumHours{} person-hours. Annotators were recruited via professional academic networks, possessed formal research backgrounds in NLP, RAG, or experimental methodology, and were compensated at standard professional rates. All participants provided explicit consent for de-identified annotation release. Institutional review determined formal IRB oversight was exempt or non-human-subjects research. Released datasets comprise labels, source text locations, rationales, and adjudicated consensus; participant identities and consent forms are withheld for privacy. Annotators operated independently without access to model outputs prior to label finalization.

\paragraph{Reproducibility and automated execution.}
All corpus-scale annotations (\S\ref{sec:tested} and \S\ref{sec:applied}) were produced via deterministic algorithms or LLM judges under fully released prompts and raw log files, enabling full re-execution. Excerpted paper quotes are checked verbatim against source line numbers by a released verifier.

\paragraph{Dual use considerations and licensing.}
Our literature audit relies exclusively on public academic publications and introduces no personal contact. Regarding dual-use security considerations: our survey discusses published threat models and defenses strictly at a conceptual level to assess experimental assumptions, releasing no attack payloads, poisoned corpora, or exploit code. All open-weight model runs utilize deterministic decoding parameters with full raw output releases.





